%%%%%%%%%%%%%%%%%%%%%%%%%%%%%% Pacotes básicos %%%%%%%%%%%%%%%%%%%%%%%%%%%%%%%
\usepackage[utf8]{inputenc}
\usepackage[T1]{fontenc}
\usepackage[brazilian,english]{babel}
\usepackage{csquotes}
\usepackage[top=3cm,bottom=2cm,right=2cm,left=3cm]{geometry}
\usepackage{cmap}

\usepackage{wasysym}
\usepackage{newunicodechar}
\newunicodechar{☺}{\smiley}
\usepackage{siunitx}

%%%%%%%%%%%%%%%%%%%%%%%%%%%%%% Fontes e matemática %%%%%%%%%%%%%%%%%%%%%%%%%%%
\usepackage{lmodern}
\usepackage{pdfpages}
\usepackage{amsmath, amssymb, amsfonts, amsthm}
\usepackage{mathrsfs}
\usepackage{lscape}
\usepackage{csquotes}
\usepackage[most]{tcolorbox}
%%%%%%%%%%%%%%%%%%%%%%%%%%%%%% Cores %%%%%%%%%%%%%%%%%%%%%%%%%%%%%%%%%%%%%%%%%
\usepackage[usenames,dvipsnames,svgnames,table]{xcolor}

%%%%%%%%%%%%%%%%%%%%%%%%%% MARCAÇÕES DE REVISÃO %%%%%%%%%%%%%%%%%%%%%%%%%%
\usepackage{xcolor}  % provavelmente já carregado em pacotes.tex

% Cores das marcações
\definecolor{marcaA}{RGB}{180, 240, 180}  % verde (Camada A: aplicada)
\definecolor{marcaB}{RGB}{255, 240, 100}  % amarelo (Camada B: sugestão)
\definecolor{marcaS}{RGB}{255, 200, 200}  % rosa (Suely)

% Reduz a borda do colorbox para não destacar demais
\setlength{\fboxsep}{1pt}

% Versão de revisão (com cores) -- ATIVA AGORA:
\newcommand{\mB}[2]{\colorbox{marcaB}{#2}\textsuperscript{\textbf{\textcolor{blue!70!black}{#1}}}}
\newcommand{\mS}[2]{\colorbox{marcaS}{#2}\textsuperscript{\textbf{\textcolor{red!70!black}{#1}}}}
\newcommand{\mA}[2]{\colorbox{marcaA}{#2}\textsuperscript{\textbf{\textcolor{green!40!black}{#1}}}}

% Versão de defesa (sem cores) -- INATIVA. Para desligar marcações antes da defesa,
% comente as três linhas acima e descomente as três abaixo:
% \newcommand{\mA}[2]{#2}
% \newcommand{\mB}[2]{#2}
% \newcommand{\mS}[2]{#2}

%%%%%%%%%%%%%%%%%%%%%%%%%%%%%% Layout %%%%%%%%%%%%%%%%%%%%%%%%%%%%%%%%%%%%%%%%
\usepackage{graphicx}
\usepackage{setspace}
\setstretch{1.5}
\usepackage{changepage}
\usepackage{indentfirst}
\setlength{\parindent}{1.5cm}

%%%%%%%%%%%%%%%%%%%%%%%%%%%%%% Cabeçalho %%%%%%%%%%%%%%%%%%%%%%%%%%%%%%%%%%%%%
\usepackage{fancyhdr}
\pagestyle{fancy}
\fancyhf{}
\fancyhead[R]{\thepage}
\renewcommand{\headrulewidth}{0pt}

\fancypagestyle{plain}{%
  \fancyhf{}
  \fancyhead[R]{\thepage}
  \renewcommand{\headrulewidth}{0pt}
}

%%%%%%%%%%%%%%%%%%%%%%%%%%%%%% Notas de rodapé %%%%%%%%%%%%%%%%%%%%%%%%%%%%%%%
\usepackage[bottom,hang]{footmisc}
\setlength{\footnotemargin}{0em}

\let\oldfootnote\footnote
\renewcommand{\footnote}[1]{%
  \oldfootnote{\fontsize{8.5}{10.5}\selectfont #1}%
}

%%%%%%%%%%%%%%%%%%%%%%%%%%%%%% Citação ABNT (ambiente) %%%%%%%%%%%%%%%%%%%%%%%
\newenvironment{citacaoabnt}{%
  \begin{singlespace}
  \begin{adjustwidth}{4cm}{0cm}
  \fontsize{9}{11}\selectfont
}{%
  \end{adjustwidth}
  \end{singlespace}
  \vspace{0.5\baselineskip}
}

%%%%%%%%%%%%%%%%%%%%%%%%%%%%%% Epígrafe %%%%%%%%%%%%%%%%%%%%%%%%%%%%%%%%%%%%%%
\newcommand{\epigrafe}[2]{%
  \vspace{0.5\baselineskip}
  \begin{flushright}
  \begin{minipage}{0.6\textwidth}
  \begin{singlespace}
  \fontsize{9}{11}\selectfont\itshape
  #1
  \end{singlespace}
  \vspace{2pt}
  \fontsize{9}{11}\selectfont --- #2
  \end{minipage}
  \end{flushright}
  \vspace{\baselineskip}
}

%%%%%%%%%%%%%%%%%%%%%%%%%%%%%% Bibliografia %%%%%%%%%%%%%%%%%%%%%%%%%%%%%%%%%%
%%%%%%%%%%%%%%%%%%%%%%%%%%%%%% Bibliografia %%%%%%%%%%%%%%%%%%%%%%%%%%%%%%%%%%
\usepackage[
  backend=biber,
  style=authoryear,
  dashed=false,      % <--- Adicione esta linha para repetir os nomes
  sortcites=true,
  sorting=nty,       
  hyperref=true,
  backref=true,
  backrefstyle=none,
  uniquename=false
]{biblatex}

\addbibresource{biblatex_style_samples/sample.bib}
\addbibresource{biblatex_style_samples/sample-abel.bib}

\DefineBibliographyStrings{brazilian}{%
  backrefpage  = {Citado na p.},
  backrefpages = {Citado nas pp.}
}

%%%%%%%%%%%%%%%%%%%%%%%%%%%%%% Tabelas 
\usepackage{booktabs}
\usepackage{longtable}
\usepackage{array}
\usepackage{multirow}
\usepackage{multicol}
\AtBeginEnvironment{longtable}{\singlespacing\fontsize{9}{11}\selectfont}
\AtBeginEnvironment{tabular}{\singlespacing\fontsize{9}{11}\selectfont}

%%%%%%%%%%%%%%%%%%%%%%%%%%%%%% Algoritmos %%%%%%%%%%%%%%%%%%%%%%%%%%%%%%%%%%%%
\usepackage{float}
\usepackage{algorithm}
\usepackage{algorithmic}
\floatname{algorithm}{Algoritmo}
\renewcommand{\listalgorithmname}{Lista de Algoritmos}

%%%%%%%%%%%%%%%%%%%%%%%%%%%%%% Código %%%%%%%%%%%%%%%%%%%%%%%%%%%%%%%%%%%%%%%%
\usepackage{tcolorbox}
\tcbuselibrary{skins,breakable,listings}

\newtcblisting{codigo}[1][]{%
  listing only,
  listing options={
    language=Python,
    basicstyle=\normalsize\ttfamily,
    breaklines=true,
    numbers=none
  },
  colback=white,
  colframe=white,
  boxrule=0pt,
  left=1.5cm
}

%%%%%%%%%%%%%%%%%%%%%%%%%%%%%% Figuras e legendas %%%%%%%%%%%%%%%%%%%%%%%%%%%%
\usepackage{caption}
\usepackage{subcaption}
\usepackage{float}

\captionsetup{
  labelfont=bf,
  textfont=normalfont
}

% BOTÃO "voltar para lista de figuras"
\usepackage{etoolbox}
\pretocmd{\listoffigures}{\phantomsection}{}{}

\DeclareCaptionFormat{figback}{#1#2#3\hfill\hyperlink{lof}{\footnotesize$\uparrow$ Lista}}
\captionsetup[figure]{format=figback}

\AtBeginDocument{
  \addtocontents{lof}{\protect\hypertarget{lof}{}}
}

%%%%%%%%%%%%%%%%%%%%%%%%%%%%%% LINKS (PROFISSIONAL) %%%%%%%%%%%%%%%%%%%%%%%%%%%
\usepackage[pdfencoding=auto]{hyperref}


\hypersetup{
  pdftitle={\titulo},
  pdfauthor={\autor},
  pdfsubject={Tese de Doutorado},

  colorlinks=true,
  linkcolor=black,
  citecolor=black,
  urlcolor=blue!50!black,

  breaklinks=true,
  bookmarksopen=true,
  bookmarksnumbered=true,
  pdfstartview=FitH,
  pdfpagemode=UseOutlines
}

%%%%%%%%%%%%%%%%%%%%%%%%%%%%%%%%%%%%%%%%%%%%%%%%%%%%%%%%%%%%%%%%%%%%%%%%%%%%%
% NÃO redefinir \frontmatter (evita erros)
%%%%%%%%%%%%%%%%%%%%%%%%%%%%%%%%%%%%%%%%%%%%%%%%%%%%%%%%%%%%%%%%%%%%%%%%%%%%%

%%%%%%%%%%%%%%%%%%%%%%%%%%%%%% Índice %%%%%%%%%%%%%%%%%%%%%%%%%%%%%%%%%%%%%%%%
\usepackage{makeidx}
\makeindex